\documentclass[11pt]{article}
\usepackage[a4paper,margin=1in]{geometry}
\usepackage[T1]{fontenc}
\usepackage{lmodern}
\usepackage{amsmath,amssymb,amsthm,mathtools,booktabs}
\usepackage[numbers,sort&compress]{natbib}
\usepackage[colorlinks=true,allcolors=blue]{hyperref}
\hypersetup{pdftitle={Energy-weighted channel recovery in a fixed-topology JT model},pdfauthor={}}
\newtheorem{theorem}{Theorem}
\newtheorem{lemma}{Lemma}
\newtheorem{proposition}{Proposition}
\newtheorem{corollary}{Corollary}
\newcommand{\tr}{\operatorname{Tr}}
\newcommand{\sech}{\operatorname{sech}}
\newcommand{\norm}[1]{\left\lVert#1\right\rVert}
\author{}
\date{29 October 2026}
\title{Energy-weighted channel recovery\\in a fixed-topology JT model}
\begin{document}
\maketitle
\begin{abstract}
We prove an energy-constrained classical-gravity recovery estimate for a
specified one-particle scalar sector of a fixed-topology Jackiw--Teitelboim
model. For the same initial gravitational Gaussian and every scalar input
with mean mode excitation at most $\nu>0$, including coherent infinite tails
and arbitrary reference systems, the reduced channel approaches its scalar
clock unitary at a sufficient rate $C^{-1/4}$. An explicit upper coefficient
is $4.531256535935\sqrt\nu$. The proof keeps the energy budget inside a
Gaussian Duhamel residual, including the direct charge-square term, before
truncating the tail. A bounded second excitation moment gives a separate
$C^{-1/2}$ estimate. The Hamiltonian, Gaussian moments, weighted charge
inequalities, and domain argument are specified here; portable code checks
the algebra and evaluates the bounds. No optimality, ADM-energy uniformity,
or global novelty claim is made.
\end{abstract}

\section{Specified model and result}
Set the AdS radius and $\hbar$ to one. We use the fixed-topology canonical
JT reduction with a free massless scalar and reflecting Dirichlet boundary
conditions. The prior canonical framework is the physical Hilbert space
and boundary generators of Penington and Witten~\citep{PeningtonWitten},
in particular their Eqs.~(41)--(45). The scalar representation follows the
positive-frequency Dirichlet branch~\citep{Higuchi2022}. On its one-particle
sector, with orthonormal modes $|n\rangle$, $n\ge0$, the charges are
\begin{align}
 K_0|n\rangle&=(n+1)|n\rangle,&
 K_+|n\rangle&=\sqrt{(n+1)(n+2)}|n+1\rangle,\notag\\
 K_-&=K_+^\dagger,& K_2&=(K_+-K_-)/(2i). \label{eq:charges}
\end{align}
The modes are $\sin[(n+1)x]e^{-i(n+1)T}/\sqrt{\pi(n+1)}$ on global
AdS$_2$, $0<x<\pi$, with metric $(-dT^2+dx^2)/\sin^2x$. The charges
preserve particle number. Restricting to this sector does not truncate
its infinite mode ladder.

Write $C=\phi_b>0$, $\chi=\log(2C)+q$, and $p=-i\partial_q$.
Our physical simultaneous-boundary time generator is
\begin{equation}
 H_C={p^2\over8C}+2Ce^{2q}+e^qK_0+{K_2^2\over8C}
 \quad\text{on }L^2(\mathbb R)\otimes\ell^2(\mathbb N_0).
 \label{eq:H}
\end{equation}
The cited boundary sum is
$2C(H_L+H_R)=p_\chi^2/4+K_2^2/4+e^\chi K_0+e^{2\chi}$,
with $J_1=-K_0$, $J_3=-K_2$; hence~\eqref{eq:H} is $H_L+H_R$
after this shift. It is not a modified decoupled Hamiltonian.
We use the positive self-adjoint realization specified in Section~\ref{sec:domains}.

The normalized initial gravity wavefunction and reference scalar clock are
\begin{equation}
 G_C(q,0)=\left({4C\over\pi}\right)^{1/4}e^{-2Cq^2},\qquad
 U_{\rm cl}(t)=e^{-i\theta(t)K_0},\quad\theta(t)=\arctan(\sinh t).
 \label{eq:initial}
\end{equation}
Let $\Phi_C(t)$ adjoin this gravity state, evolve with $H_C$, and trace
gravity, and let $\Phi_{\rm cl}(t)=\operatorname{Ad}_{U_{\rm cl}(t)}$.
Define the unhalved energy-constrained distance
\begin{equation}
 d_\nu(C,T)=\sup_{\substack{0\le t\le T,\ R,\ \rho_{mR}\ge0\\
 \tr\rho_{mR}=1,\ \tr(n\rho_m)\le\nu}}
 \norm{[(\Phi_C(t)-\Phi_{\rm cl}(t))\otimes\operatorname{Id}_R]
 (\rho_{mR})}_1,\quad n=K_0-1. \label{eq:distance}
\end{equation}
Here $R$ ranges over reference systems. No diagonality, finite second
moment, or finite ADM mean is assumed.

Put
\begin{align}
 M(v)&=e^{8v}-2e^{2v}(1+6v+8v^2)+1+8v+12v^2,\notag\\
 F(v)&=e^{2v}-2e^{v/2}+1,\qquad g_C=2C\sqrt{M(1/(8C))}.\label{eq:MF}
\end{align}
For integer $n\ge0$ define the nonnegative coefficient polynomials
\begin{equation}
 u(n)={(n+1)(n+2)^2(n+3)\over16},\quad
 d(n)={(n+1)^4\over4},\quad
 \ell(n)={n^2(n^2-1)\over16},\label{eq:polynomials}
\end{equation}
and, for $N\ge1$,
\begin{equation}
 B_N(\nu)=\sqrt{{3\over4}+{u(N)-3/4\over N}\nu}
 +\sqrt{{1\over4}+{d(N)-1/4\over N}\nu}
 +\sqrt{{\ell(N)\over N}\nu}.\label{eq:B}
\end{equation}

\begin{theorem}\label{thm:main}
For every integer $N\ge1$, $C>0$, $\nu\ge0$, and $0\le T\le1$,
\begin{align}
 d_\nu(C,T)\le\min\Bigg\{2,\;&
 2T\left[g_C+\sqrt{1+(N+2)\nu}\sqrt{F(1/(8C))}
             +{B_N(\nu)\over8C}\right]\notag\\
 &+4\sqrt{{\nu\over N+1}}\Bigg\}.\label{eq:finite}
\end{align}
For each fixed $\nu>0$, choosing
$N=\lfloor(4\sqrt2/3)\sqrt C\rfloor$ gives
\begin{equation}
 d_\nu(C,1)\le {14\,2^{3/4}\sqrt3\over9}\sqrt\nu\,C^{-1/4}
                    +O_\nu(C^{-1/2}).\label{eq:rate}
\end{equation}
For $\nu=0$, the input is the ground mode and
$d_0(C,T)\le2T[g_C+\sqrt{F(1/(8C))}+1/(8C)]=O(C^{-1/2})$.
\end{theorem}

\section{Normalized Gaussian reference and its residual}
The reference gravitational path and stability solution are
\begin{equation}
 q_c=-\log\cosh t,\quad p_c=-4C\tanh t,\quad
 z(t)=1-t\tanh t+i\tanh t. \label{eq:path}
\end{equation}
They obey $\dot q_c=p_c/(4C)$, $\dot p_c=-4Ce^{2q_c}$,
$z''+2\sech^2(t)z=0$, $z(0)=1$, and $z'(0)=i$.
With $\xi=q-q_c(t)$ and $S_C(t)=2C[t-2\tanh t]$, the Gaussian
\begin{equation}
 G_C(q,t)=\left({4C\over\pi}\right)^{1/4}z(t)^{-1/2}
 \exp\left(i\left[S_C+p_c\xi+2C{z'\over z}\xi^2\right]\right)
 \label{eq:gaussian}
\end{equation}
uses the square root continuous from $z(0)=1$. The Wronskian gives
$\operatorname{Im}(z'/z)=|z|^{-2}$; thus it is normalized and its position
variance is $v_C(t)=|z(t)|^2/(8C)$. Substitution shows that it solves
\begin{equation}
 H_g^{(2)}(t)={p^2\over8C}+2Ce^{2q_c(t)}(1+2\xi+2\xi^2).
 \label{eq:Hquad}
\end{equation}
The phase in~\eqref{eq:gaussian} includes the constant term of this
quadratic Taylor Hamiltonian. Writing $y=\tanh t$ gives
$|z|^2-1=y[(1+t^2)y-2t]\le0$ on $[0,1]$, since $y\le t$ and
$1+t^2\le2$.

For a finite-sequence scalar/reference vector $f$, set
$f(t)=(U_{\rm cl}(t)\otimes I)f$. The product
$\Psi_{\rm app}=G_C(t)\otimes f(t)$ has exact residual
\begin{align}
 (H_C-i\partial_t)\Psi_{\rm app}={}&
 2Ce^{2q_c}(e^{2\xi}-1-2\xi-2\xi^2)G_C\otimes f(t)\notag\\
 &+e^{q_c}(e^\xi-1)G_C\otimes K_0f(t)
 +G_C\otimes {K_2^2f(t)\over8C}.\label{eq:residual}
\end{align}
All three terms are retained. In particular the direct charge-square term
is not replaced by a scalar expectation value.

For a centered normal variable $X$ of variance $v$,
$\mathbb E e^{aX}=e^{a^2v/2}$. Differentiating this identity gives
\begin{align}
 \mathbb E[e^{2X}(1+2X+2X^2)]&=e^{2v}(1+6v+8v^2),\notag\\
 \mathbb E(1+2X+2X^2)^2&=1+8v+12v^2.
\end{align}
Consequently the two squared residual norms are exactly
\begin{equation}
 \mathbb E(e^{2X}-1-2X-2X^2)^2=M(v),\qquad
 \mathbb E(e^X-1)^2=F(v).\label{eq:momentderivation}
\end{equation}
Both functions increase for $v\ge0$. All nonconstant coefficients of
$F$ are positive; those of $M$ vanish below degree three and at degree
$k\ge3$ equal $2^k[4^k-2(2k^2+k+1)]/k!>0$. Their leading terms are
$M(v)=(80/3)v^3+O(v^4)$ and $F(v)=v+O(v^2)$.
Since $e^{q_c},e^{2q_c}\le1$, Duhamel's identity bounds the vector error,
uniformly for $t\le T\le1$, by
\begin{equation}
 T\left[g_C\norm f+\sqrt{F(1/(8C))}\norm{K_0f}
                  +{\sup_{s\le T}\norm{K_2^2f(s)}\over8C}\right].
 \label{eq:duhamel}
\end{equation}
This compares exact and reference evolutions from the same initial vector.
The scalar phase preserves mode probabilities and $\norm{K_0f}$.

\section{Energy-weighted charge bound and coherent tail}
Purify an input of~\eqref{eq:distance}, enlarging its reference if needed.
Let $P_N=\sum_{n=0}^N|n\rangle\langle n|$ and $f=(P_N\otimes I)\psi$.
Then $\norm f\le1$ and, with reference-valued coefficients $f_n$,
\begin{equation}
 \sum_{n=0}^N n\norm{f_n}^2\le\nu,\qquad
 \norm{K_0f}^2\le1+(N+2)\nu.\label{eq:weightedK0}
\end{equation}
The second inequality uses $(n+1)^2=1+n(n+2)$ and $n\le N$.
Squaring the \emph{infinite} charge in~\eqref{eq:charges} gives
\begin{align}
 K_2^2|n\rangle={}&-{(n+2)\sqrt{(n+1)(n+3)}\over4}|n+2\rangle
               +{(n+1)^2\over2}|n\rangle\notag\\
 &-{n\sqrt{(n-1)(n+1)}\over4}|n-2\rangle,\label{eq:K2square}
\end{align}
where the downward term is zero at $n=0,1$. Its three squared coefficients
are~\eqref{eq:polynomials}. Input through mode $N$ can produce output
$N+2$; squaring a truncated charge would omit this contribution.

\begin{lemma}\label{lem:charge}
If $\norm f\le1$, $f\in\operatorname{Ran}(P_N\otimes I)$, and
$\sum n\norm{f_n}^2\le\nu$, then $\norm{K_2^2f}\le B_N(\nu)$.
\end{lemma}
\begin{proof}
For integer $0\le n\le N$,
\begin{align}
 u(n)&\le\tfrac34+{n\over N}[u(N)-\tfrac34],&
 d(n)&\le\tfrac14+{n\over N}[d(N)-\tfrac14],&
 \ell(n)&\le{n\over N}\ell(N).\label{eq:chords}
\end{align}
The first two quotients after subtracting constants are
\begin{equation}
 {u(n)-3/4\over n}={n^3\over16}+{n^2\over2}+{23n\over16}+{7\over4},
 \qquad {d(n)-1/4\over n}={n^3\over4}+n^2+{3n\over2}+1,
\end{equation}
increasing for $n\ge0$. The last is $(n^3-n)/16$, increasing at integers
$n\ge1$, and $\ell(0)=0$. Each of the three shifts in~\eqref{eq:K2square}
has squared norm equal to its coefficient polynomial averaged against
$\norm{f_n}^2$. Apply~\eqref{eq:chords}, then the triangle inequality
between shifts. This gives~\eqref{eq:B} for reference-valued coefficients.
\end{proof}

\begin{proof}[Proof of Theorem~\ref{thm:main}]
Equation~\eqref{eq:duhamel}, \eqref{eq:weightedK0}, and
Lemma~\ref{lem:charge} bound the projected vector error by the bracket
in~\eqref{eq:finite} times $T$. These bounds hold at all reference times
because phases preserve mode probabilities. Both Stinespring vectors have
norm $\norm f\le1$, so their reduced subnormalized density matrices
differ in trace norm by at most twice their vector distance.

The omitted probability obeys
$p=\norm{(1-P_N)\psi}^2\le\nu/(N+1)$. In the span of the retained and
omitted components, the pure-state density minus its subnormalized
projection has trace norm
\begin{equation}
 \norm{|\psi\rangle\langle\psi|-|f\rangle\langle f|}_1
       =\sqrt{4p-3p^2}\le2\sqrt p.\label{eq:tail}
\end{equation}
Partial tracing preserves the upper bound. Both channels contract Hermitian
trace norm; their difference contributes at most $4\sqrt p$ on the omitted
part. This proves~\eqref{eq:finite}, uniformly over references, without
differentiating an infinite-tail vector.

For fixed $\nu>0$, $B_N(\nu)=\sqrt\nu N^{3/2}+O_\nu(N^{1/2})$.
With $N=\lfloor a\sqrt C\rfloor$, the leading coefficient is
\begin{equation}
 \sqrt\nu\left(\sqrt{a/2}+{a^{3/2}\over4}+{4\over\sqrt a}\right)C^{-1/4}.
\end{equation}
The gravity term is only $2g_C=\sqrt{5/6}\,C^{-1/2}+O(C^{-3/2})$.
The leading coefficient is stationary at
$1/\sqrt2+3a/4-4/a=0$, with positive solution $a=4\sqrt2/3$.
This proves~\eqref{eq:rate}. For $\nu=0$ the input class is exactly
the ground mode and $\norm{K_0|0\rangle}=\norm{K_2^2|0\rangle}=1$.
\end{proof}

\section{Bounded second moment: a separate input class}
If $\tr(n^2\rho_m)\le\mu$, the projected vector satisfies
$\mathbb E n\le\sqrt\mu$ and $\mathbb E n^j\le N^{j-2}\mu$
for $j=2,3,4$. Thus $\norm{K_0f}\le1+\sqrt\mu$ and
\begin{align}
 B_N^{(2)}(\mu)={}&\sqrt{\tfrac34+\tfrac74\sqrt\mu+
                    (N^2/16+N/2+23/16)\mu}\notag\\
 &+\sqrt{\tfrac14+\sqrt\mu+(N^2/4+N+3/2)\mu}
                    +\tfrac14N\sqrt\mu\label{eq:B2}
\end{align}
bounds $\norm{K_2^2f}$. Expand $u,d,\ell$ and use the same shift
triangle inequality to obtain this expression. The tail probability
is now at most $\mu/(N+1)^2$.

\begin{corollary}
For the distance with the stronger constraint $\tr(n^2\rho_m)\le\mu$,
make the following replacements in~\eqref{eq:finite}:
\[
 \sqrt{1+(N+2)\nu}\mapsto1+\sqrt\mu,\quad
 B_N(\nu)\mapsto B_N^{(2)}(\mu),\quad
 4\sqrt{\nu/(N+1)}\mapsto {4\sqrt\mu\over N+1}.
\]
Taking $N=\lfloor4\sqrt C\rfloor$ gives
\begin{equation}
 d_{\mu;2}(C,1)\le
 \left[\sqrt{5/6}+{1+\sqrt\mu\over\sqrt2}+2\sqrt\mu\right]C^{-1/2}
       +O_\mu(C^{-1}).\label{eq:second}
\end{equation}
\end{corollary}
No fourth moment of the full input is required: only its finite projection
enters Duhamel, and bounded exact channels act on the tail.

\section{Domains, quantitative consequences, and limits}\label{sec:domains}
The realization of~\eqref{eq:H} is associated with the closed nonnegative form
\begin{equation}
 \mathfrak h_C[\psi]={\norm{p\psi}^2+\norm{K_2\psi}^2\over8C}
     +2C\norm{e^q\psi}^2+\norm{e^{q/2}K_0^{1/2}\psi}^2.\label{eq:form}
\end{equation}
Its domain is the intersection of the four indicated operator domains.
The Jacobi charge has its standard self-adjoint closure; its linearly
growing coefficients satisfy the divergent reciprocal-coefficient
criterion. The finite sum of closed nonnegative forms on this dense
intersection is closed. Smooth compact geometry functions times finite
scalar sequences form a common form core. To see this, first use smooth
scalar cutoffs $f(n/N)$: the commutator $[K_2,f(n/N)]$ is uniformly bounded
and supported in a moving tail, hence tends strongly to zero on every
vector. These cutoffs converge in the $K_2$ graph norm and commute with
$K_0$ and geometry operators. Smooth compact geometry cutoffs then have
vanishing derivative commutators; smoothing on compact support completes
the approximation. This does not assume sharp scalar projections converge
in an unbounded graph norm.

For the finite scalar inputs used in Duhamel, applying the differential
expression to the Gaussian reference produces finite sums of Gaussian
polynomials times exponentials and finite scalar sequences. Every derivative
and exponential multiple in~\eqref{eq:form} is square integrable.
Integration by parts against the common form core therefore puts these
trial vectors in the operator domain, with action~\eqref{eq:residual}.
Their time dependence is continuously differentiable on the compact
interval, so Duhamel's identity is valid. General energy-constrained inputs
enter through the trace-norm tail step; they need not be in
$\operatorname{Dom}H_C$.

Table~\ref{tab:bounds} gives outward-rounded 150-decimal interval bounds.
The earlier worst-code projection argument substitutes $N+1$ and
$(N+1)(N+2)$ for the weighted norms; $N=\lfloor(8\nu C)^{1/3}\rfloor$
then gives $3\sqrt2\nu^{1/3}C^{-1/6}+O(C^{-1/3})$.
The new bound retains the mean budget inside the residual, rather than
re-optimizing the same worst-code estimate. At $N\asymp\sqrt C$ its
constrained vector error vanishes although the worst-code operator
estimate need not. For fixed $\nu$, the sufficient small-error resource
scaling becomes $C=O(\nu^2\delta^{-4})$, compared with
$O(\nu^2\delta^{-6})$ from that earlier argument. These are sufficient
asymptotics, not necessary thresholds.

\begin{table}[htbp]
\centering
\caption{Unhalved channel-distance upper bounds at $\nu=1$, $T=1$;
every displayed value is rounded upward.}\label{tab:bounds}
\begin{tabular}{@{}rrr@{}}\toprule
$C$ & Energy-weighted bound & Worst-code projection bound\\\midrule
$10^6$ & $0.144226$ & $0.435331$\\
$10^{12}$ & $0.004532171$ & $0.042528$\\
$10^{18}$ & $0.000143292$ & $0.004244$\\
$10^{24}$ & $0.000004531258$ & $0.000424275$\\\bottomrule
\end{tabular}
\end{table}

The portable script checks exact charge and Gaussian-moment identities,
the stability equation and Wronskian, and charge columns retaining two
upper-edge raising steps. Its 180 coherent matter/reference samples support
the inequalities; arbitrary-reference coverage follows from the proof.
As a hostile tail check,
$\sqrt{1-1/N}|0\rangle+\sqrt{1/N}|N\rangle$ has mean excitation one,
but $\norm{K_0\psi}^2=N+3$ and projection error $\sqrt{4/N-3/N^2}$.
A mean-as-norm substitution or probability-only coherent-tail bound fails.

The constraint is scalar global-mode excitation, not total ADM energy;
finite mean $K_0$ alone need not imply finite $K_2$ form energy.
This is leading scalar-clock recovery, distinct from a quantum-corrected
finite-code channel expansion. No optimal exponent, arbitrary-time,
topology-changing, general-Fock, boundary-instrument, or four-dimensional
GR conclusion follows. The computation evaluates explicit bounds and
finite coefficient checks; it does not simulate the full JT propagator.

\section*{Provenance and disclosure}
The canonical JT framework cited above is prior work. The Gaussian
comparison was inherited from an internal September 2026 calculation;
the energy-weighted improvement was derived on 28 October 2026 and is
packaged here on 29 October 2026. Historical source dates, this delivery,
and demonstrable public disclosure are distinct: a private commit does
not establish public priority. No global novelty certification is claimed.

Earlier work was Fable 5-led; July 2026 onward results were predominantly
produced by GPT models, as reported by the operator. DeepSeek and multiple
GPT models also contributed cross-review and manuscript preparation.
This is reported chronology, not a per-result authorship allocation.
Human names and affiliations are blank, and human expert review has not
occurred. The supplied statement is preserved verbatim:
\begin{quote}
``I genuinely do not know how authorship, intellectual credit, or ownership should be attributed for these results.''
\end{quote}
\bibliographystyle{unsrtnat}
\bibliography{references}
\end{document}
